\chapter{Deployment and Operations}
\label{ch:operations}

\epigraph{``Get busy living, or get busy dying.''}{Andy Dufresne,
\emph{The Shawshank Redemption} (1994)}

Operations is unglamorous compounding. Health checks that mean something, drains
that actually drain, alerts that fire before a user notices rather than after.
None of it is clever, and all of it is what makes three in the morning
uneventful.

The good news for MCP specifically is that the stateless core deleted most of
the hard parts of this chapter before it was written. If you deployed a 2025-era
server you spent real effort on session affinity, a session store, and draining
state during a deploy. All three of those problems are simply gone.

\section{What statelessness bought you}

Worth listing explicitly, because if you deployed a 2025-era MCP server you spent
real effort on all of these.

\begin{table}[htbp]
\centering
\small
\begin{tabularx}{\textwidth}{@{}lXX@{}}
\toprule
\thead{Concern} & \thead{Handshake era} & \thead{2026-07-28} \\
\midrule
Load balancing   & Sticky sessions, session affinity & Round robin \\
Session store    & Redis, or affinity you regret & None \\
Scaling out      & Drain sessions, migrate state & Add a replica \\
Deploys          & Sessions break or must migrate & Requests retry \\
Serverless       & Awkward at best & First-class \\
Crash recovery   & Rebuild session state & Lose one request \\
\bottomrule
\end{tabularx}
\caption{The right-hand column is why this chapter is shorter than it would have
been a year ago.}
\label{tab:statelessness-ops}
\end{table}

\keyidea{If you find yourself needing sticky sessions for an MCP server, you have
per-request state somewhere it should not be. Find it. It is usually a task
store, a signing secret, or a cache with a per-process lifetime.}

\section{Three topologies}

\subsection{stdio sidecar}

One server process per host, launched as a subprocess. No network.

\textbf{Good for:} local tools, filesystem access, developer machines, anything
where the user's own credentials are the right credentials.

\textbf{Operational profile:} no ports, no TLS, no auth infrastructure. The
process boundary is the isolation boundary. Cold start is 44.6\,ms, which
Chapter~\ref{ch:optimizing} says to amortise with a pool.

\textbf{Watch for:} the server inheriting more environment than it should. A
subprocess gets the parent's environment by default, including credentials for
things it has no business touching.

\begin{py}
# sketch: condensed for the page
self._proc = subprocess.Popen(
    command,
    stdin=subprocess.PIPE, stdout=subprocess.PIPE,
    stderr=subprocess.PIPE if capture_stderr else None,
    env={**os.environ, **(env or {})},
    ...
)
\end{py}

Meridian merges the parent environment because it is a book. A hardened launcher
passes an explicit allowlist instead.

\subsection{Containerised Streamable HTTP}

The default for anything shared, and the whole deployment is a container behind
a load balancer:

\begin{py}
FROM python:3.13-slim
WORKDIR /app
COPY meridian/ ./meridian/
ENV PYTHONUNBUFFERED=1
EXPOSE 8931
CMD ["python3", "-m", "meridian.serve", "risk", "--http", "8931"]
\end{py}

\texttt{PYTHONUNBUFFERED=1} matters more than it looks: buffered stderr means
your logs arrive in 4\,KB chunks, which during an incident means they arrive
after you needed them.

Behind it, a plain round-robin load balancer. No affinity configuration. That is
the entire deployment story now.

\subsection{Serverless}

\textbf{Good for:} spiky traffic, low steady-state volume, and anything where you
would rather not run a fleet.

\textbf{The catch, and it is the whole catch:} cold start. Your first request
pays interpreter boot plus imports plus whatever runs at module scope. Meridian's
\texttt{servers/data.py} builds 240 accounts, 1,440 transactions, and 120 filings
at import time, which is fine at 40\,ms and would not be at four seconds.

Two rules for serverless MCP servers:

\begin{itemize}[leftmargin=1.6em, itemsep=3pt]
\item \textbf{Nothing expensive at module scope.} Lazy-load anything heavy, so a
cold container answers \texttt{tools/list} without loading your scoring model.
\item \textbf{Tasks need external storage.} Obviously, but people forget: a
container that answered \texttt{tools/call} will not be the one that answers
\texttt{tasks/get}.
\end{itemize}

\section{Health checks that mean something}

Most MCP health checks are wrong, and in an interesting way. They check that the
process is listening, which tells you almost nothing.

A useful health check exercises the protocol:

\begin{py}
def healthy(server) -> tuple[bool, dict]:
    """Health means `server/discover` answers correctly, not that a port is open.

    Answering discovery proves routing, capability derivation, and
    serialisation all work, which is three subsystems for one request.
    """
    probe = {
        "jsonrpc": "2.0", "id": "health", "method": "server/discover",
        "params": {"_meta": {
            KEY_PROTOCOL_VERSION: PROTOCOL_VERSION,
            KEY_CLIENT_CAPABILITIES: {},
        }},
    }
    response = server.handle(probe)
    if response is None or "error" in response:
        reason = (response or {}).get("error", {}).get("message", "no response")
        return False, {"reason": reason}

    result = response["result"]
    ok = (PROTOCOL_VERSION in result.get("supportedVersions", [])
          and bool(result.get("capabilities")))
    return ok, {
        "capabilities": sorted(result.get("capabilities", {})),
        "supportedVersions": result.get("supportedVersions", []),
    }
\end{py}

Three levels, and use all three:

\begin{table}[htbp]
\centering
\small
\begin{tabularx}{\textwidth}{@{}lXl@{}}
\toprule
\thead{Check} & \thead{Proves} & \thead{Frequency} \\
\midrule
Liveness: process responds
  & It has not deadlocked & seconds \\
Readiness: \texttt{server/discover} returns capabilities
  & Routing and serialisation work & seconds \\
Deep: one real tool call against a known fixture
  & Dependencies are reachable & minutes \\
\bottomrule
\end{tabularx}
\caption{The deep check is the one that catches a server whose database went
away, and it is the one most people skip.}
\label{tab:health-checks}
\end{table}

\begin{notebox}{Do not use \texttt{ping}}
It was removed in \texttt{2026-07-28}. If your health check sends it, a modern
server answers \texttt{-32601} and your orchestrator concludes the server is
unhealthy and restarts it, forever.

The dual-era bridge answers \texttt{ping} for exactly this reason, which is a compatibility shim, not a recommendation.
\end{notebox}

\section{Graceful drain}

Statelessness removes most of the difficulty here, because there is no session
state to migrate and no affinity to preserve. There is still a right order, and
step three is the one specific to MCP.

\begin{enumerate}[leftmargin=1.6em, itemsep=3pt]
\item \textbf{Fail readiness.} The load balancer stops sending new requests.
\item \textbf{Finish in-flight requests.} They are short. Give them a bounded
window.
\item \textbf{Close \texttt{subscriptions/listen} streams gracefully.} This is
the MCP-specific step and it is the one people miss.
\item \textbf{Exit.}
\end{enumerate}

Step three deserves attention. A subscription stream that just dies tells the
client nothing. A stream that closes with the response to its originating
request tells it ``this ended on purpose'':

\begin{py}
def close_gracefully(self) -> None:
    """End the subscription with a proper JSON-RPC response.

    A stream that just dies tells the client nothing. A stream that closes
    with the response to the original `subscriptions/listen` request tells
    it "this ended on purpose, do not reconnect in a panic". The difference
    matters at three in the morning.
    """
    self._send_raw(jsonrpc.result_response(
        self.id, {"_meta": {KEY_SUBSCRIPTION_ID: self.id}}))
    self.closed = True
\end{py}

Without it, every client treats your rolling deploy as an outage and reconnects
immediately, all at once, to the replicas that are still up.

On stdio, drain means honouring EOF on stdin:

\begin{py}
def close(self, timeout: float = 5.0) -> None:
    try:
        self._proc.stdin.close()
    except Exception:
        pass
    try:
        self._proc.wait(timeout=timeout)
    except subprocess.TimeoutExpired:
        # Escalate only if the server ignored the EOF on stdin.
        self._proc.terminate()
\end{py}

\section{Rolling out a capability change}

The awkward case: you are adding a tool, and for a while some replicas have it
and some do not.

\begin{dangerbox}{The split-catalogue window}
Replica A has the new tool. Replica B does not. A client fetches
\texttt{tools/list} from A, caches it for five minutes, then calls the new tool
and hits B, which returns \texttt{-32602}.

Round-robin load balancing plus client-side caching makes this the \emph{normal}
case during a deploy, not an edge case.
\end{dangerbox}

Three mitigations, and use the first two always:

\textbf{Deploy the handler before the declaration.} Ship a release where the tool
works but is not listed. Then ship a release that lists it. Now no replica ever
advertises something it cannot do.

\textbf{Keep TTLs short enough that the window closes.} Meridian's
\texttt{tools/list} TTL is five minutes, so a deploy is fully visible within five
minutes of finishing.

\textbf{Emit \texttt{listChanged} when the rollout completes.} Clients with an
open subscription refetch immediately instead of waiting out the TTL.

Removal runs in reverse: stop listing it, wait for the TTL plus a margin, then
delete the handler. And if it is a public server, apply the deprecation window
from Chapter~\ref{ch:architecture} to your own surface. Your callers deserve the
same twelve months the specification gives you.

\section{Multi-tenancy without sessions}

Tenancy used to ride on the session. There is no session, so it rides on the
credential, which is better.

\begin{py}
def visible_tools(self, ctx: RequestContext) -> list[Tool]:
    """Override to filter by scope. The set may vary by authorization,
    never by connection."""
\end{py}

Four rules:

\textbf{Tenant comes from the token, never from an argument.} A tenant id in
\texttt{arguments} is a suggestion from a model that read attacker-controlled
text. A tenant id in a validated token claim is a fact.

\textbf{Fold tenancy into every cache key.} Meridian folds the auth context in
unconditionally, public or private, precisely so a mislabelled response cannot
leak across tenants.

\textbf{Bind MRTR state to the principal.} Chapter~\ref{ch:mrtr} covers it. Across
tenants it is the difference between an approval and an incident.

\textbf{Put the tenant in \texttt{baggage} for tracing, and nothing else.}
Sliceable traces, no sensitive data crossing trust boundaries.

\section{Alerting}

Chapter~\ref{ch:measuring} covered the protocol alerts. These are the
operational ones, and each has an MCP-specific flavour.

\begin{table}[htbp]
\centering
\small
\begin{tabularx}{\textwidth}{@{}lXX@{}}
\toprule
\thead{Alert} & \thead{Symptom} & \thead{Usual cause} \\
\midrule
Tool-call loop
  & Same tool, same arguments, many times in one task
  & Unhelpful error message; the model cannot self-correct \\
Subscription storm
  & \texttt{subscriptions/listen} count spikes
  & A deploy dropped streams without graceful closure \\
Timeout spike on one method
  & p99 up on \texttt{tools/call}, flat elsewhere
  & One tool's dependency is degraded \\
Cache stampede
  & Request spike at a TTL boundary
  & Synchronised TTLs; add jitter \\
Task store growth
  & Store size climbing without bound
  & Expiry sweep not running \\
\bottomrule
\end{tabularx}
\caption{Five alerts. The first is the one that costs money silently.}
\label{tab:ops-alerts}
\end{table}

The tool-call loop deserves elaboration because it is uniquely agentic. A model
calls a tool, gets an unhelpful error, retries identically, gets the same error,
and repeats until the step budget stops it. Every cycle is a full model turn.
Nothing is down, nothing errors at the infrastructure level, and your bill goes
up.

The fix is Chapter~\ref{ch:tools}: error messages a model can act on. The alert
is a detector for having failed to do that.

Cache stampede is worth a sentence too. A thousand clients that started together
have synchronised TTLs, so they all expire together and all refetch together.
Jitter the TTL by a few percent per client and the herd disperses.

\section{Capacity planning}

The useful thing about MCP servers: they are ordinary request-response services
and your existing capacity model applies.

The unusual thing: your traffic shape is driven by \emph{loop iterations}, not by
user actions. One user request becomes three model turns becomes three tool
calls, sometimes fanned out in parallel bursts.

So plan for the burst. From Chapter~\ref{ch:measuring}, work forward:

\begin{plain}
tasks/second                    x  round trips/task
                                =  requests/second (mean)

requests/second (mean)          x  fan-out width
                                =  concurrent requests at the burst
\end{plain}

Meridian at 100 tasks/second: 3 round trips each is 300 requests/second mean,
and a fan-out of 3 means bursts of up to 900 concurrent. Provisioning for the
mean gives you a service that falls over during every parallel fan-out, which is
to say during every task.

\subsection{Little's Law, which sizes your pools}

That arithmetic has a name, and knowing it turns capacity planning from
guesswork into a division.

Little's Law says that for any stable system, the average number of items inside
it equals the arrival rate multiplied by the average time each spends inside:

\begin{plain}
L = lambda * W

L       concurrent requests in flight
lambda  arrival rate, requests per second
W       average time in the system, seconds
\end{plain}

It holds regardless of the arrival distribution, the service-time distribution,
or the scheduling order. The only requirement is that the system is stable, which
is to say arrivals are not outrunning departures. That generality is what makes
it worth memorising.

Now use it. Meridian's compliance server handles 300 requests per second with a
mean service time of 8 milliseconds:

\begin{plain}
L = 300 * 0.008 = 2.4 concurrent requests
\end{plain}

Two and a half. A thread pool of 8 looks generous, and this is precisely how
people arrive at a pool that collapses under a fan-out, because the mean is the
wrong input. Substitute the burst instead: 900 concurrent arrivals at 8
milliseconds needs $900 \times 0.008 = 7.2$, and if a downstream dependency
degrades and service time goes to 200 milliseconds, the same arrival rate needs
$900 \times 0.2 = 180$.

That last calculation is the important one, and it is why the ``slow dependency''
row in the load-test table exists. Concurrency requirements scale linearly with
service time, so a dependency that gets 25 times slower needs 25 times the
in-flight capacity to sustain the same throughput. A pool sized for the healthy
case starts queueing, queueing adds latency, added latency raises $W$, and a
higher $W$ needs a bigger $L$ than you have. That loop is what a congestion
collapse looks like from the inside.

Run the same division for every bounded resource in the path: worker threads,
the HTTP connection pool from Chapter~\ref{ch:optimizing}, database connections,
and the file descriptors all three consume. The one you did not size is the one
that fails first.

\keyidea{Size pools with $L = \lambda W$, using your burst arrival rate and your
\emph{degraded} service time. Sizing for the mean of both is how a service that
looks comfortably provisioned falls over the first time something downstream
gets slow.}

\section{Load testing to failure}

Load testing an agent system exercises something different from load testing an
API, because the traffic shape is generated by loop iterations rather than by
users pressing buttons.

\begin{table}[htbp]
\centering
\small
\begin{tabularx}{\textwidth}{@{}lXX@{}}
\toprule
\thead{Test} & \thead{What it finds} & \thead{Why MCP-specific} \\
\midrule
Sustained load
  & Steady-state capacity & Baseline \\
Burst load
  & Fan-out behaviour & Real traffic is bursty by construction \\
Cold cache
  & Load with no client caching & Your first minute after a deploy \\
Split catalogue
  & Half the replicas on the new version & The deploy window above \\
Slow dependency
  & Behaviour when a downstream is degraded & Timeouts consume step budgets \\
\bottomrule
\end{tabularx}
\caption{The last row is the interesting one: in an agent system a slow
dependency does not just add latency, it consumes the caller's step budget and
produces no information.}
\label{tab:load-tests}
\end{table}

\subsection{What breaks first}

From running Meridian's compliance server under load, in the order it happened:

\begin{enumerate}[leftmargin=1.6em, itemsep=3pt]
\item \textbf{Thread pool exhaustion under fan-out.} Bursts are three times mean
concurrency and the pool was sized for the mean.
\item \textbf{The audit log became the bottleneck.} Every call writes a record
synchronously, so the log's write path caps your throughput. Batch it, or write
it asynchronously with a durable queue in front.
\item \textbf{SSE streams accumulated.} Clients that vanish without closing leave
streams open. Keep-alives detect it; without them, file descriptors leak until
the process stops accepting connections.
\end{enumerate}

None of those are protocol problems. All three are consequences of the traffic
shape the protocol produces.

\section{Configuration and secrets}

\begin{table}[htbp]
\centering
\small
\begin{tabularx}{\textwidth}{@{}lXl@{}}
\toprule
\thead{Setting} & \thead{Notes} & \thead{Scope} \\
\midrule
MRTR signing secret & Rotatable, verify old and new during rotation & \textbf{fleet} \\
Task store URL      & Shared and durable & fleet \\
TTL values          & Tune per resource, jitter per client & fleet \\
Auth issuer, audience & Per environment & fleet \\
Page size           & Balance round trips against response size & per server \\
Extension toggles   & Only if implemented & per server \\
\bottomrule
\end{tabularx}
\caption{The bold row is the one that produces intermittent, replica-count-
proportional failures when it is wrong.}
\label{tab:config}
\end{table}

Rotating the MRTR secret needs care, because an in-flight elicitation may be
sealed with the old key. Accept both during the rotation window, sign only with
the new one, and make the window longer than the state TTL. Otherwise every
approval in flight at rotation time fails, and the failure looks exactly like
tampering.

\section{Meridian in production}

\begin{plain}
                     internet
                        |
                 [ API gateway ]     TLS, OAuth, rate limits,
                        |            routes on Mcp-Method / Mcp-Name
        +---------------+---------------+
        |               |               |
   [ risk x3 ]   [ compliance x2 ]  [ fraud x3 ]   [ marketdata x2 ]
        |               |               |               |
        +---------------+---------------+---------------+
                        |
              [ shared task store ]     durable, expiring
              [ signing secret    ]     from the secret manager
\end{plain}

Every replica count above one is possible because there is nothing to make
sticky. The gateway routes on headers without parsing bodies, caches
\texttt{public} results from market data, and enforces per-method rate limits.

\begin{measurebox}{What the gateway cache is worth}
\begin{tabular}{@{}lrr@{}}
\toprule
& \thead{Per-task setup} & \thead{Requests, 20 tasks} \\
\midrule
No client caching   & 274.1\,ms & 164 \\
Honouring \texttt{ttlMs} & 2.5\,ms & 8 \\
\bottomrule
\end{tabular}

\vspace{6pt}
\footnotesize
Client-side. A shared gateway cache in front of \texttt{marketdata}, whose
results are \texttt{public}, removes most of the remainder for every client at
once.
\end{measurebox}

\section{The runbook}

A runbook is what you want written down before you need it. The test of one is
whether somebody who did not build the system can follow it at four in the
morning.

\begin{table}[htbp]
\centering
\small
\begin{tabularx}{\textwidth}{@{}lX@{}}
\toprule
\thead{Symptom} & \thead{First thing to check} \\
\midrule
Intermittent failures, roughly $1-1/N$
  & Per-process state: signing secret, task store, cache \\
Clients cannot connect at all
  & Era mismatch. Probe with \texttt{server/discover} by hand \\
Everything times out on stdio
  & Something wrote to \texttt{stdout} \\
Latency up, server metrics flat
  & Iterations per task. It is a round-trip problem \\
Cost up, latency flat
  & Catalogue tokens. Somebody added tools \\
Cache hit rate collapsed
  & TTLs changed, or \texttt{tools/list} ordering went unstable \\
Streaming stopped working through the proxy
  & \texttt{X-Accel-Buffering: no} \\
Approvals failing since the deploy
  & Signing secret rotation window too short \\
\bottomrule
\end{tabularx}
\caption{Eight symptoms, eight first moves. Every one of them appears somewhere
earlier in this book.}
\label{tab:runbook}
\end{table}

\section{What to remember}

\begin{itemize}[leftmargin=1.6em, itemsep=3pt]
\item Statelessness deleted sticky sessions, session stores, and session
migration. If you still need affinity, find the per-process state.
\item Health checks should call \texttt{server/discover}, not open a socket. Do
not use \texttt{ping}; it was removed.
\item Drain in order, and close subscription streams gracefully or every deploy
looks like an outage.
\item Deploy the handler before the declaration; remove the declaration before
the handler.
\item Tenancy comes from the token, folds into every cache key, and binds MRTR
state.
\item Plan capacity for the fan-out burst, not the mean.
\item Alert on tool-call loops. They cost money while every dashboard stays
green.
\item Jitter your TTLs.
\item Rotate the MRTR signing secret with an overlap window longer than the state
TTL.
\end{itemize}

That closes Part~IV. You can build it, measure it, make it fast, and run it.
Part~V is about the people trying to break it.
